% Transcription of the user's vkv_draft0.pdf (received 2026-10-23 as a chat attachment; the original file was not
% on disk). Text and formulas are kept verbatim, including the author's wording; only typesetting is redone.
\documentclass[11pt]{article}
\usepackage[a4paper,margin=2.2cm]{geometry}
\usepackage{amsmath,amssymb,amsthm,booktabs,enumitem,xcolor}
\usepackage[linesnumbered,ruled,vlined]{algorithm2e}
\usepackage[hidelinks]{hyperref}
\usepackage[most]{tcolorbox}
\setlength{\parskip}{4pt}
\setlength{\parindent}{0pt}
\newcommand{\eps}{\varepsilon}
\newtheorem{lemma}{Lemma}
\newtheorem{proposition}[lemma]{Proposition}
\newtheorem{theorem}[lemma]{Theorem}
\theoremstyle{definition}
\newtheorem{assumption}{Assumption}
\renewcommand{\theassumption}{A\arabic{assumption}}
\theoremstyle{remark}
\newtheorem{remark}{Remark}

\title{\textbf{Vkv}\\[2pt]\large Truncated Horvitz--Thompson estimation with certified fetching for the attention denominator}
\author{VerifiedKV algorithm explainer}
\date{23 October 2026}

\begin{document}
\maketitle

At prefill, Vkv keeps each tail token with probability $r_i$ (KVzip scores) and moves the rest to CPU. At every attention
decision it does three things, all computed from the query sketch and $r$, never from which tokens were drawn. First, it
\emph{truncates} the tail tokens whose Horvitz--Thompson upweighting is too risky: if sampled, they get weight 1 instead of
$1/r$. The set is chosen by an MSE criterion. Second, it bounds the resulting error with a Bernstein certificate that treats
the known truncation bias and the random fluctuation separately. Third, if the bound exceeds the tolerance $\eps_0$, it
fetches the heaviest tokens by sketch weight, and picks how many to fetch in a single prefix-sum sweep. With exact weights
in the plan, every decision then satisfies $\mathbb P(|\hat B - B| > \eps_0 B) \le \delta$. The effect of sketch error is
left to Part III.

\part*{Part I \quad Theory}

\section{Setting}

Fix one (layer, KV head) and one decision, i.e.\ one query row. The context has keys $k_i$, and the query $q$ (scaled by
$1/\sqrt d$) gives weights $w_i = e^{\langle q, k_i\rangle - m}$. The quantity to estimate is the softmax denominator
\[ B = \sum_i w_i + w^x, \]
where $w^x$ is the exact weight of the question and generated keys.

\textbf{Prefill.} Each context token has a keep probability $r_i = \min(1, c\tau_i)$ from KVzip scores $\tau_i$. Sinks and
the local window get $r_i = 1$. Write $H = \{r_i = 1\}$ for the always-resident set and $L = \{r_i < 1\}$ for the tail.
Independent draws $Z_i \sim \mathrm{Bernoulli}(r_i)$ decide which tail tokens stay on the GPU; call these the sample
$S = \{i \in L : Z_i = 1\}$. The full cache is offloaded to the CPU.

\textbf{Sketch.} A compact sketch of all keys gives an approximate weight $\hat w_i$ for every tail token, sampled or not.
In the analysis below we take the plan's weights to be exact ($\hat w_i = w_i$). The effect of sketch error is discussed
in Part III.

\begin{assumption}[Draw independence]\label{a1}
The $Z_i$ are independent, and independent of everything a plan may use: the query, the sketch, and $r$.
\end{assumption}

This holds in the question pass when draws are independent across layers. At decode it holds along a fixed reference
prefix.

\section{The base estimator and why plans must ignore the draw}

The Horvitz--Thompson (HT) estimator weights each sampled tail token by $1/r_i$:
\[ \hat B_{\mathrm{HT}} = \sum_H w_i + w^x + \sum_{i \in L} \frac{Z_i}{r_i} w_i, \qquad \mathbb E \hat B_{\mathrm{HT}} = B,
\qquad \mathrm{Var} = \sum_L w_i^2 \frac{1 - r_i}{r_i}. \]

\textbf{Design rule.} Every set used below (truncated, fetched) must be chosen from $(\hat w, r)$ only. If you instead pick
sampled tokens \emph{because} they were sampled and give them weight 1, you discard the mass they represented. On a
synthetic head this biased the estimate by $-25\%$ to $-41\%$.

\section{Truncation: trading a little bias for a lot of variance}

\subsection{One token}

Give a tail token weight 1 instead of $1/r$ if it is sampled. Its contribution changes as follows:
\begin{center}
\begin{tabular}{lccc}\toprule
 & weight if sampled & bias & variance \\ \midrule
HT & $1/r$ & $0$ & $w^2(1-r)/r$ \\
Truncated & $1$ & $-w(1-r)$ & $w^2 r(1-r)$ \\ \bottomrule
\end{tabular}
\end{center}
So truncation adds bias $b = w(1-r)$, always downward, and removes variance $w^2(1-r)/r - w^2 r(1-r) = w^2(1-r)^2(1+r)/r$.

\subsection{Many tokens: why the per-token test is wrong}

$\mathrm{MSE} = (\text{total bias})^2 + \text{total variance}$. Variances of independent tokens add. Biases add
\emph{before} squaring, and they all point the same way. If $T$ is the truncated set and $\beta_T = \sum_{i \in T}
w_i(1 - r_i)$ its total bias, the bias term is $\beta_T^2$. Keep in mind the square of sum vs sum of square.

\subsection{The MSE rule}

Adding token $j$ to $T$ changes the MSE by
\[ \Delta\mathrm{MSE} = \underbrace{2\beta_T w_j(1 - r_j) + w_j^2(1 - r_j)^2}_{(\beta_T + b_j)^2 - \beta_T^2}
 - w_j^2 (1 - r_j)^2 \frac{1 + r_j}{r_j}. \]
Dividing by $w_j(1 - r_j)$ gives
\begin{tcolorbox}[colback=green!3,colframe=green!40!black,boxrule=0.4pt,arc=2pt]
\centering truncating $j$ lowers the MSE $\iff J_j := \dfrac{w_j(1 - r_j)}{r_j} > 2\beta_T$.
\end{tcolorbox}
$J_j$ is token $j$'s \emph{jump}: how far HT overshoots when $j$ is sampled. Go through tokens in decreasing order of jump.
Jumps fall and $\beta_T$ rises along that order, so the MSE first decreases and then increases, and the best set $T^\star$
is the prefix that stops at the first $j$ with $J_j \le 2\beta_T$. Large-jump tokens (heavy weight, tiny $r$) are truncated
first.

\begin{lemma}[Truncation caps the jumps]\label{lem:cap}
Every tail token left as HT has $J_j \le 2\beta_{T^\star}$.
\end{lemma}
\begin{proof}
The rule stops at the first token whose jump is at most $2\beta_{T^\star}$, and all later tokens in the order have jumps no
larger.
\end{proof}

This is what makes the certificate in \S\ref{sec:cert} tight. Without truncation, a single token with $r = 10^{-3}$ can have a
jump $1000\times$ its weight, and any Bernstein-type bound pays for that even though it happens with probability $10^{-3}$.

\begin{remark}
\textbf{interesting thing i found:} Shrinking extreme HT weights to trade bias for variance is known in survey statistics as
\emph{weight trimming}. What is specific here is choosing the trimmed set from a query-time sketch, before the draw is seen,
and using the trimming to bound the jumps so that a closed-form certificate becomes tight.
\end{remark}

\section{Fetching}

Fetching a token makes it exact. For this to stay unbiased, \emph{every} fetched token must get weight 1, not only the ones
outside the sample. Tokens outside $S$ are read from CPU; tokens already sampled lose their $1/r$ upweighting. (If sampled
members kept $1/r$, a fetched token would contribute $w + w(1 - r)$ on average.) A fetched token then contributes no error.
The transfer cost is the number of fetched tokens that were not sampled.

\textbf{Order.} Fetch tail tokens in decreasing order of $\hat w$, as in needle fetch. Let $F_k$ be the top $k$. Fetching
removes a token's variance and, if it was truncated, also its bias.

\section{The error after fetching \texorpdfstring{$k$}{k} tokens}

After fetching the top $k$ tokens, every tail token falls into one of three groups. Each group contributes to the error
$\hat B_k - B$ in a different way.
\begin{itemize}
\item \textbf{Fetched} ($F_k$): counted with its exact weight. Error 0.
\item \textbf{Truncated, not fetched} ($T^\star \setminus F_k$): counted with weight 1 if sampled, missing if dropped.
\item \textbf{HT, not fetched} (the rest of $L$): counted with weight $1/r$ if sampled, missing if dropped.
\end{itemize}
Tokens in $H$ and the question/generated keys are exact and contribute no error.

\textbf{One truncated token.} If it was sampled (probability $r$), we count $w$, which is exactly right: error 0. If it was
dropped (probability $1 - r$), we count nothing: error $-w$. So its error is $w(Z - 1)$. On average this is $-w(1 - r)$
(the bias basically), a \emph{known} shortfall we can compute in advance. Around that average, the error moves up by
$w(1 - r)$ when the token is sampled, and down by $wr$ when it is dropped.

\textbf{One HT token.} If it was sampled, we count $w/r$: error $+w(1 - r)/r$ (its jump). If it was dropped, we count
nothing: error $-w$. So its error is $w(Z/r - 1)$. On average this is exactly 0 (this is why we call this part unbiased):
the occasional overshoot cancels the frequent miss.

\begin{center}
\begin{tabular}{lccccc}\toprule
 & if sampled & if dropped & mean & variance & max move up / down \\ \midrule
Truncated & $0$ & $-w$ & $-w(1-r)$ & $w^2 r(1-r)$ & $w(1-r)$ \ / \ $wr$ \\
HT & $+w(1-r)/r$ & $-w$ & $0$ & $w^2(1-r)/r$ & $w(1-r)/r$ \ / \ $w$ \\ \bottomrule
\end{tabular}
\end{center}

\textbf{Adding up.} Each unfetched token's error is its mean plus a random part with mean 0 (we do this so we can later apply
the bernstein, which only works on mean zero vector sum). Summing over all unfetched tokens splits the total error into a
known part and a random part:
\begin{equation}\label{eq:split}
\hat B_k - B = \underbrace{-\beta(k)}_{\text{known bias}} + \underbrace{Y(k)}_{\text{random part}}, \qquad
\beta(k) = \sum_{T^\star \setminus F_k} w_i (1 - r_i).
\end{equation}
\begin{itemize}
\item $\beta(k)$ is the sum of the truncated tokens' shortfalls. It does not depend on the draw, and it always pushes the
estimate down.
\item $Y(k)$ is the sum of the random parts of all unfetched tokens. They are independent (each token was kept
independently) and each has expectation value zero with non zero variance. Their total variance is
\[ V(k) = \sum_{T^\star \setminus F_k} w_i^2 r_i(1 - r_i) + \sum_{L \setminus (T^\star \cup F_k)} w_i^2 \frac{1 - r_i}{r_i}, \]
and no single one moves away from its mean by more than $b(k)$, the largest entry of the last column over unfetched tokens.
\end{itemize}

\textbf{Why this split helps.} The known bias can simply be added to the error budget. Only $Y(k)$ needs a probability
bound, and for that we need just $V(k)$ and $b(k)$ (next section). Fetching moves a token out of the residual groups, which
removes its terms from $\beta$, $V$ and $b$. That is how fetching lowers the error bound. By Lemma~\ref{lem:cap}, the upward
moves of the remaining HT tokens are at most $2\beta_{T^\star}$, which keeps $b(k)$ small.

\section{The certificate}\label{sec:cert}

\begin{lemma}[Bernstein]
If $Y$ is a sum of independent, zero-mean terms with $|Y_i| \le b$ and $\sum \mathrm{Var}(Y_i) = V$, then for $t > 0$,
$\mathbb P(|Y| \ge t) \le 2\exp\bigl(-t^2 / (2(V + bt/3))\bigr)$.
\end{lemma}

Set the right-hand side equal to $\delta$ and solve. With $L_\delta = \ln(2/\delta)$,
\[ t(k) = \frac{b(k) L_\delta}{3} + \sqrt{\Bigl(\frac{b(k) L_\delta}{3}\Bigr)^2 + 2V(k)L_\delta}, \qquad
\boxed{\eps(k) = \frac{\beta(k) + t(k)}{B}} \]
On the event $|Y(k)| \le t(k)$, which has probability at least $1 - \delta$, \eqref{eq:split} gives
$|\hat B_k - B| \le \beta(k) + t(k) = \eps(k) B$.

\section{Decision and fetch size in one rule}

\begin{proposition}[Monotone along the order]\label{prop:mono}
$\beta(k)$, $V(k)$ and $b(k)$ never increase in $k$, so neither does $\eps(k)$.
\end{proposition}
\begin{proof}
Each fetch removes a non-negative term from $\beta$ and from $V$, and a maximum over a smaller set cannot grow. $t$ is
increasing in $V$ and $b$.
\end{proof}

\textbf{Rule.} If $\eps(0) \le \eps_0$, do not fetch. Otherwise fetch $F_{k^\star}$, where $k^\star$ is the smallest $k$
with $\eps(k) \le \eps_0$. By Proposition~\ref{prop:mono}, $k^\star$ is the first index that passes in one sweep. It always
exists, because fetching the whole tail gives $\beta = V = b = 0$.

\section{The guarantee}

\begin{theorem}[Vkv guarantee]\label{thm:vkv}
Under Assumption~\ref{a1}, let Vkv fetch $F_{k^\star}$ with $k^\star$ the smallest $k$ such that $\eps(k) \le \eps_0$
(no fetch if $k = 0$ passes). Then for every decision,
\[ \mathbb P\bigl(|\hat B_{k^\star} - B| > \eps_0 B\bigr) \le \delta. \]
Consequently the expected fraction of failing decisions is at most $\delta$, whatever the correlation between decisions.
\end{theorem}
\begin{proof}
$T^\star$, the fetch order and $\eps(k)$ depend only on the weights and $r$, so by Assumption~\ref{a1} the chosen $k^\star$
is independent of the draw. The certificate of \S\ref{sec:cert} then applies at $k^\star$. The last claim is linearity of
expectation.
\end{proof}

\begin{remark}[Marginal targets]
The theorem certifies each decision separately. If transfer per forward turns out too high (one fetch set serves all query
rows in a forward), replace the per-decision test by a budget on the \emph{average} failure bound across a forward. The
expected failing fraction stays bounded by linearity, and the policy can give up on a few expensive diffuse decisions.
\end{remark}

\part*{Part II \quad The full algorithm}

Algorithm~\ref{alg:prefill} runs once per context. Algorithm~\ref{alg:decision} runs for every decision, vectorised over
query rows. Every quantity in lines 1--18 of Algorithm~\ref{alg:decision} is computed from $(\hat w, r)$ alone; the draw $Z$
is used only to decide what to transfer and how to bias attention (lines 19--22).

\begin{algorithm}[H]
\caption{Vkv prefill (once per context)}\label{alg:prefill}
\KwIn{context keys and values; kept fraction $\kappa$}
Compute KVzip scores $\tau_i$\;
Choose $c$ so that $\sum_i \min(1, c\tau_i) = \kappa n$; set $r_i \gets \min(1, c\tau_i)$\;
$r_i \gets 1$ on sinks and the local window; $H \gets \{r_i = 1\}$, $L \gets \{r_i < 1\}$\;
\ForEach{(layer, KV head)}{
  Draw $Z_i \sim \mathrm{Bernoulli}(r_i)$ independently for $i \in L$ (independent stream per (layer, KV head))\;
  Keep $H \cup \{i \in L : Z_i = 1\}$ on the GPU\;
}
Offload the full cache to CPU; build the key sketch; store $r_i$ for all $i \in L$\;
\end{algorithm}

\begin{algorithm}[H]
\caption{Vkv decision, truncation and fetch (one query row)}\label{alg:decision}
\KwIn{sketch weights $\hat w_i$ for $i \in L$; exact $w_i$ on $H$ and $w^x$; $r$; tolerance $\eps_0$; failure
probability $\delta$}
\KwOut{fetch set $F$; attention bias}
$\tilde B \gets \sum_H w_i + w^x + \sum_L \hat w_i$; \ $L_\delta \gets \ln(2/\delta)$\;
\tcp{(a) truncation set by the MSE rule}
$J_i \gets \hat w_i (1 - r_i)/r_i$ for $i \in L$; sort $L$ by $J$ in decreasing order\;
$T^\star \gets \emptyset$; \ $\hat\beta \gets 0$\;
\ForEach{$i$ in that order}{
  \lIf{$J_i \le 2\hat\beta$}{break}
  $T^\star \gets T^\star \cup \{i\}$; \ $\hat\beta \gets \hat\beta + \hat w_i(1 - r_i)$\;
}
\tcp{(b) per-token bias, variance and maximal move}
\For{$i \in L$}{
  \eIf{$i \in T^\star$}{
    $b_i \gets \hat w_i(1 - r_i)$; \ $v_i \gets \hat w_i^2 r_i (1 - r_i)$; \ $m_i \gets \max(\hat w_i(1 - r_i),\, \hat w_i r_i)$\;
  }{
    $b_i \gets 0$; \ $v_i \gets \hat w_i^2 (1 - r_i)/r_i$; \ $m_i \gets \max(\hat w_i(1 - r_i)/r_i,\, \hat w_i)$\;
  }
}
\tcp{(c) one sweep in needle order, $k = 0, 1, \dots, |L|$}
Sort $L$ by $\hat w$ in decreasing order; let $F_k$ be the first $k$ tokens\;
$\hat\beta(k) \gets \sum_{L \setminus F_k} b_i$; \ $\hat V(k) \gets \sum_{L \setminus F_k} v_i$; \
$\hat b(k) \gets \max_{L \setminus F_k} m_i$ \tcp*{prefix sums, suffix max}
$t(k) \gets \dfrac{\hat b(k) L_\delta}{3} + \sqrt{\Bigl(\dfrac{\hat b(k) L_\delta}{3}\Bigr)^2 + 2\hat V(k) L_\delta}$\;
$\hat\eps(k) \gets \bigl(\hat\beta(k) + t(k)\bigr) / \tilde B$\;
\tcp{(d) decision}
\eIf{$\hat\eps(0) \le \eps_0$}{
  $F \gets \emptyset$ \tcp*{no fetch}
}{
  $k^\star \gets \min\{k : \hat\eps(k) \le \eps_0\}$; \ $F \gets F_{k^\star}$; \ fetch $F \setminus S$ from CPU\;
}
\tcp{(e) attention bias}
Bias 0 on $H \cup F$ and on sampled tokens in $T^\star \setminus F$\;
Bias $\log(1/r_i)$ on sampled tokens in $L \setminus (T^\star \cup F)$\;
Bias $-\infty$ on all other tail tokens\;
\end{algorithm}

The attention bias in (e) implements the estimator
\[ \hat B = \sum_H w_i + w^x + \sum_F w_i + \sum_{S \cap T^\star \setminus F} w_i + \sum_{S \setminus (T^\star \cup F)}
\frac{w_i}{r_i}, \]
where $S = \{i \in L : Z_i = 1\}$. By Proposition~\ref{prop:mono}, $\hat\eps(k)$ is non-increasing, so $k^\star$ is the first
index that passes in the sweep. With exact weights ($\hat w = w$, so $\tilde B = B$), Theorem~\ref{thm:vkv} gives
$\mathbb P(|\hat B - B| > \eps_0 B) \le \delta$.

\part*{Part III \quad Implementation note}

\section{What the code runs (differs from the theory only where stated)}

\textbf{Tight certificate.} The bias always points down, so bound the two sides separately and pass if
\[ P_{\mathrm{fail}}(k) = e^{-u_+^2 / (2(V + b^+ u_+/3))} + e^{-u_-^2 / (2(V + b^- u_-/3))} \le \delta, \qquad
u_\pm = \eps_0 \tilde B \pm \beta(k), \]
with $b^\pm$ the separate up and down maxima from the table in \S5, and the second term set to 1 if $u_- \le 0$. Use the same
sweep and take the first passing $k$.

\begin{tcolorbox}[colback=blue!3,colframe=blue!50!black,boxrule=0.4pt,arc=2pt,title=Why this is safe]
If the simple check passes at $k$, then $t \le u_- \le u_+$, so each term is at most $\delta/2$. The tight check therefore
passes wherever the simple one does: the guarantee carries over, and $k_{\mathrm{practice}} \le k_{\mathrm{theory}}$. What is
lost is \emph{monotonicity}. As $\beta$ shrinks, $u_+$ shrinks too, so the over-estimation term can grow: $P_{\mathrm{fail}}$
can pass at $k$ and fail at $k + 1$ (112 of 20k rows in simulation). The first passing index is still the minimal passing
one, but the tight check cannot be binary-searched, and ``fetching more is always safe'' does not hold for it. Both
properties hold for the simple check only.
\end{tcolorbox}

\textbf{Sketch error (theory deferred).} With real sketches, $\hat\beta$, $\hat V$, $\hat b$ and $\tilde B$ are estimates,
so the certificate is approximate. First run with \emph{true} weights in place of $\hat w$ (the simulation has them) to
separate algorithm behaviour from sketch error. A simple correction to try later: inflate $V$ by $\gamma^2$, $b$ and $\beta$ by
$\gamma$, and divide $\tilde B$ by $\gamma$, with $\gamma$ calibrated so these hold on most decisions.

\textbf{Numerics and cost.}
\begin{itemize}
\item Two sorts, a few cumsums and one reverse cummax per row. A top-$K$ partial sort (e.g.\ $K = 4096$) suffices; treat the
rest as HT.
\item Compute $\hat w^2(1 - r)/r$ in fp32, since it overflows fp16 when $r$ is tiny.
\item $T^\star$ lives in the per-row attention bias, so truncation costs no transfer.
\item Keep $T^\star$ frozen after fetching. Re-truncating could only lower the bound; freezing may over-fetch slightly but
stays valid.
\item If Bernstein is still loose, try exact Chernoff at a few ladder points rather than in the sweep. In the synthetic test it
was only about $1.3\times$ tighter.
\end{itemize}

\section{What to log, and what ``logical'' looks like}

Per decision, log: phase, layer, head; $\eps(0)$; decision; $k^\star$; transfers $|F \setminus S|$; $|T^\star|$;
$\beta/\tilde B$ at $k^\star$; $b^+$ before and after truncation. Using the full cache, also log the true error
$e = |\hat B_{k^\star} - B|/B$, and the error the no-fetch estimator would have had.

\begin{center}\small
\begin{tabular}{p{0.38\textwidth}p{0.55\textwidth}}\toprule
Check & Expected if Vkv is behaving \\ \midrule
Fetch rate vs $\eps_0$, $\delta$, kept fraction & Rises as each decreases; monotone curves \\
$\mathbb P(e > \eps_0)$ over all decisions & $\le \delta$; with true weights, a violation is a bug \\
$k^\star$ vs oracle $k$ (smallest $k$ with true error $\le \eps_0$) & Ratio $\ge 1$; its size is the bound's slack \\
Tight vs simple certificate & Tight never fetches more per row; log rows where $P_{\mathrm{fail}}$ passes at $k$ and fails at
$k + 1$ \\
$b^+$ before vs after truncation & After $\le 2\beta$ (Lemma~\ref{lem:cap}); before can be orders of magnitude larger \\
Truncation on vs off ($T^\star = \emptyset$) & Truncation cuts fetches at matched failure rate \\
Transfer per forward (union over rows) & Compare with needle fetch at matched memory; if far higher, switch to the averaged
budget (Remark 2) \\
$\beta/\tilde B$ at $k^\star$ & Small (a few \%); watch whether it compounds downstream \\ \bottomrule
\end{tabular}
\end{center}

\textbf{Red flags.}
\begin{itemize}
\item Failure rate above $\delta$ with true weights: a sampled set leaked into $T^\star$ or the order, or sampled members of
$F$ kept weight $1/r$.
\item $T^\star$ equal to the whole tail: the bias is being summed over sampled tokens only. It must include unsampled truncated
tokens.
\item Fetch rate near 100\%: fp16 overflow in the variance, or a sketch that overestimates weights.
\end{itemize}

\end{document}
